%%%%%%%%%%%%%%%%%%%%%%%%%%%%%%%%%%%%%%%%%%%%%%%%%%%%%%%%%%%%%%%%%%%%%
%% sec_implementation.tex
%% Section 3: Implementation
%% RMG Periodic RT-TDDFT — Paper 2
%%%%%%%%%%%%%%%%%%%%%%%%%%%%%%%%%%%%%%%%%%%%%%%%%%%%%%%%%%%%%%%%%%%%%

\section{Implementation}
\label{sec:implementation}

The implementation of the periodic and spin-polarized RT-TDDFT
presented in this work is built on the existing finite-system
RT-TDDFT module of \texttt{RMG},\cite{jakowski-rmg-tddft-2025}
which is described in detail in the accompanying paper.
Here we focus on the extensions required for periodic boundary
conditions and spin-polarized calculations, the key approximations
and their physical justification, and the computational organization
of the velocity gauge propagation.

%--------------------------------------------------------------------
\subsection{Algorithm}
\label{sec:algorithm}
%--------------------------------------------------------------------

The computational scheme for periodic RT-TDDFT in \texttt{RMG}
is summarized in Algorithm~\ref{alg:tddft-periodic}.
It follows the same density matrix propagation framework as the
finite-system implementation~\cite{jakowski-rmg-tddft-2025} but
with three key modifications: (1) the perturbation is applied as
a momentum kick via the vector potential $\mathbf{A}$ rather than
an electric field dipole kick; (2) the current operator matrix
replaces the dipole matrix as the primary observable; and (3)
the Hamiltonian matrix is evaluated at each k-point over the
Brillouin zone.

\begin{algorithm}[h!]
\caption{Periodic RT-TDDFT Algorithm (VECTOR\_POT mode)}
\label{alg:tddft-periodic}
\begin{algorithmic}[1]
\scriptsize

\item[{\bf /* Ground state DFT calculations: */}]
\State Generate Bloch orbitals $|\psi_{n,\mathbf{k}}\rangle$,
       Kohn-Sham Hamiltonian $\vb{H}^\sigma_{0,\mathbf{k}}$
       for each k-point $\mathbf{k}$ and spin $\sigma$
\State Select active space: $N_{\mathrm{occ}}$ occupied +
       $N_{\mathrm{virt}}$ virtual orbitals
\item[]

\item[{\bf /* Build current operator matrices (once, before time loop): */}]
\State For each k-point $\mathbf{k}$ and spin $\sigma$:
\State \quad Build paramagnetic current matrix:
\item[] \quad $J^\alpha_{ab} \leftarrow
         \langle\psi_a|(-\imath\nabla_\alpha + k_\alpha)|\psi_b\rangle$
       \quad [VecPHmatrix()]
\State \quad Add nonlocal PP correction:
\item[] \quad $J^\alpha_{ab} \leftarrow J^\alpha_{ab} +
         \langle\psi_a|[\hat{r}_\alpha, \Vnl]|\psi_b\rangle$
       \quad [CurrentNlpp()]
\State \quad Add $\mathbf{A}(t_0)\cdot\mathbf{J}$ to $\vb{H}_{-1,\mathbf{k}}$
       to encode the perturbation kick at $t=0$
\item[]

\item[{\bf /* Initialize density matrix: */}]
\State $\vb{P}^0_{\mathbf{k}} \leftarrow \mathrm{diag}(f_{n,\mathbf{k}})$
       \quad (occupation numbers on diagonal)
\State $\vb{H}_{-1,\mathbf{k}} \leftarrow \vb{H}_{0,\mathbf{k}}$
       (extrapolation seed)
\item[]

\item[{\bf /* Time loop: */}]
\For{$j = 1$ to $N_{\mathrm{steps}}$}
  \State $t \leftarrow t_0 + j \cdot \Delta t$
  \State \textbf{Extrapolate:}
         $\vb{H}_{1,\mathbf{k}} \leftarrow 2\vb{H}_{0,\mathbf{k}}
          - \vb{H}_{-1,\mathbf{k}}$ \quad [extrapolate\_Hmatrix()]
  \item[]
  \item[{\bf /* Self-consistent Hamiltonian loop: */}]
  \While{$\|\vb{H}_{1}^{\mathrm{new}} - \vb{H}_{1}^{\mathrm{old}}\|_\infty
          > H_{\mathrm{thrs}}$}
    \State \textbf{Magnus operator:}
           $\vb{\Omega}_\mathbf{k} \leftarrow
            -\imath\cdot\tfrac{1}{2}(\vb{H}_{0,\mathbf{k}}
            + \vb{H}_{1,\mathbf{k}})\Delta t$
           \quad [magnus()]
    \State \textbf{Propagate density matrix:}
           $\vb{P}^1_\mathbf{k} \leftarrow
            \vb{P}^0_\mathbf{k} + [\vb{\Omega}_\mathbf{k}, \vb{P}^0_\mathbf{k}]
            + \ldots$ \quad [eldyn\_ort() / commutp()]
    \State \textbf{Update density:}
           $\rho^\sigma(\mathbf{r}) \leftarrow \rho^\sigma_{\mathrm{gs}}(\mathbf{r})
            + \sum_\mathbf{k} w_\mathbf{k}
              \sum_{n} P^1_{nn,\mathbf{k}} |\psi_{n,\mathbf{k}}(\mathbf{r})|^2$
           \quad [GetNewRho\_rmgtddft()]
    \State \textbf{Update potentials:}
           $V_{\mathrm{H}}[\rho]$, $V^\sigma_{\mathrm{xc}}[\rho^\uparrow,\rho^\downarrow]$
    \State \textbf{Update Hamiltonian:}
           $\vb{H}_{1,\mathbf{k}}^{\mathrm{new}} \leftarrow
            \langle\psi_a|\hat{H}[\rho]|\psi_b\rangle$
           \quad [HmatrixUpdate()]
  \EndWhile
  \item[]
  \State \textbf{Extract current:}
         $J_\alpha(t) \leftarrow
          \mathrm{Re}\sum_{\mathbf{k}} w_\mathbf{k}\,
          \mathrm{Tr}[\vb{P}^1_\mathbf{k}\cdot\vb{J}^\alpha_\mathbf{k}]$
  \State (Optional) \textbf{Berry phase polarization:}
         $P_{\mathrm{BP},\alpha}(t) \leftarrow
          \mathrm{Tr}[\vb{P}^1_\mathbf{k}\cdot\vb{X}^\alpha_\mathbf{k}]$
  \State \textbf{Advance:}
         $\vb{P}^0 \leftarrow \vb{P}^1$;\;
         $\vb{H}_{-1} \leftarrow \vb{H}_{0}$;\;
         $\vb{H}_{0} \leftarrow \vb{H}_{1}$
\EndFor
\item[]
\item[{\bf /* Postprocessing: */}]
\State Fourier transform $J_\alpha(t)$ for each Cartesian direction $\alpha$
\State Compute $\sigma_{\alpha\beta}(\omega)$,
       $\varepsilon_{\alpha\beta}(\omega)$

\end{algorithmic}
\end{algorithm}


%--------------------------------------------------------------------
\subsection{Key Approximations and Their Justification}
\label{sec:approximations}
%--------------------------------------------------------------------

The implementation described in Algorithm~\ref{alg:tddft-periodic}
embodies several approximations with respect to the full theory
presented in Section~\ref{sec:theory}.
We discuss each in turn.

\paragraph{Perturbative momentum kick.}
The full velocity gauge Hamiltonian (eq.~\ref{eq:tdks-velocity})
contains the vector potential $\mathbf{A}(t)$ at every time step,
requiring the $\mathbf{A}(t)\cdot\hat{\mathbf{p}}$ term to be
re-evaluated throughout the propagation.
In the present implementation, we instead apply $\mathbf{A}(t)$
as a delta-function kick at $t=0$ only: the momentum kick is encoded
in the initial Hamiltonian $\vb{H}_{-1}$ by adding the
$\mathbf{A}(t_0)\cdot\vb{J}^\alpha$ matrix (step 6 in
Algorithm~\ref{alg:tddft-periodic}), after which $\mathbf{A}=0$
for all $t>0$.
This approach is equivalent to first-order perturbation theory in
$|\mathbf{A}_0|$, and it provides the full linear optical response
of the system through the free oscillations of the current density
$\mathbf{J}(t)$ following the kick.
The linear response regime is valid as long as the kick amplitude
is small enough that the system does not leave the linear regime,
which is straightforwardly verified by checking that the response
is proportional to $|\mathbf{A}_0|$.

\paragraph{Diamagnetic current term.}
The full current density contains both paramagnetic and diamagnetic
contributions (Section~\ref{sec:current}).
In the present implementation, the diamagnetic term
$\mathbf{J}_{\mathrm{dia}} \propto \mathbf{A}(t)|\psi|^2$
is identically zero for all $t>0$ because $\mathbf{A}(t>0)=0$
in the perturbative kick scheme.
This is consistent with the linear response regime: the diamagnetic
contribution to the linear response function is implicitly accounted
for by the sum rule, and the paramagnetic current alone gives the
correct optical conductivity.\cite{yabana-bertsch-1996,Mattiat-Luber-2022}
Simulations with a continuously applied driving field
(e.g., a CW laser or a Gaussian pulse) would require retaining
$\mathbf{A}(t)$ throughout the propagation and are beyond the scope
of the present work.

\paragraph{First-order Magnus expansion.}
As in the finite-system implementation,\cite{jakowski-rmg-tddft-2025}
the Magnus operator is truncated to first order:
\begin{equation}
  \vb{\Omega} \approx \vb{\Omega}_1
  = -\imath \cdot \tfrac{1}{2}(\vb{H}_0 + \vb{H}_1)\Delta t.
  \label{eq:magnus-1st}
\end{equation}
The second-order correction $\vb{\Omega}_2 = \frac{1}{12}\Delta t^2
[\vb{H}_0, \vb{H}_1]$ is not included.
The accuracy and stability of this scheme for the time steps used
here was validated in paper~1 through energy conservation tests,
and we demonstrate analogous stability for the periodic case in
Section~\ref{sec:results}.

\paragraph{Norm-conserving pseudopotentials only.}
The current implementation supports only norm-conserving
pseudopotentials (NCPP).\cite{kleinman-bylander-1982,hamann-1989}
Ultrasoft pseudopotentials (USPP)\cite{vanderbilt-1990} introduce
a non-trivial overlap matrix $\mathbf{S} \neq \mathbf{I}$ into
the density matrix propagation (the nonorthogonal basis path),
which requires the generalized propagation scheme described in
Appendix A of paper~1.
This extension is planned for future work.

\paragraph{BCH propagator for complex Hamiltonians.}
For k-points away from the $\Gamma$ point, the Hamiltonian matrix
$\vb{H}_\mathbf{k}$ is complex Hermitian.
In this case, the BCH commutator expansion propagator (Ieldyn=1)
is used throughout; the diagonalization-based propagator (Ieldyn=2)
is not currently supported for complex matrices.
The BCH expansion is unconditionally stable and has been shown to
give results equivalent to exact diagonalization within the
convergence threshold,\cite{jakowski-rmg-tddft-2025} and we find
no practical limitation from this choice.

%--------------------------------------------------------------------
\subsection{Software Optimization and Parallelization}
\label{sec:optimization}
%--------------------------------------------------------------------

\myblue{[Content outline -- adapt from paper 1 Section 3.2 with ]} 
additions for the periodic case:
- Same MPI + OpenMP + GPU parallelization as paper 1
- k-point parallelization: each MPI group handles a subset of k-points
  (pct.kpsub\_comm); current is summed over k-points with weights
  (MPI\_Allreduce over kpsub\_comm, lines 828-829 in RmgTddft.cpp)
- Spin parallelization: separate MPI communicator groups per spin channel
  (pct.spin\_comm); spin channels communicate only for XC update
  (MPI\_Allreduce over spin\_comm, line 767)
- The density matrix P and Hamiltonian H are distributed using
  ScaLAPACK layout (Sp object); size is numst x numst where
  numst = N\_occ + N\_virt
- For VECTOR\_POT mode: momentum matrix J\_alpha is complex even at
  Gamma point (stored as complex<double> in Pxmatrix\_cpu etc.)
- The BCH commutator $[\Omega, P]$ dominates $O(N^3)$; same optimization
  as paper 1 (real/imaginary decomposition, BLAS ZGEMM/DGEMM)
- Timing data for benchmark systems: \myblue{[add when available]}%]}

\myblue{[Note: include a brief timing comparison table analogous to
Table 1 in paper 1, but for one or more periodic benchmark systems
(e.g., Si supercells of increasing size) to demonstrate scaling.]}
